% FILE: CSE-17_TermFinal.tex
\documentclass{article}
\usepackage{amsmath}
\usepackage{amssymb}
\begin{document}

%%%%%%%%%%%%%%%%%%%%%%%%%%%%%%%%%%%%%%%%%%%%%%%%%%
% QUESTION_START
%%%%%%%%%%%%%%%%%%%%%%%%%%%%%%%%%%%%%%%%%%%%%%%%%%
% id=cse17-final-q5ai
% discipline=CSE
% year=2017
% exam_type=Term Final
% section=B
% ct_number=UNKNOWN
% question_number=5ai
% topic=Indefinite Integration
% subtopic=General
% difficulty=Medium
% length=Medium
% question_type=New
% frequency=1
% appearances=
% OCR_STATUS=VERIFIED
\section*{Term Final (Sec B) - Q5(a)(i)}
Question:
Integrate the following:
\[
\int \frac{x^2+1}{x^4+1} \, dx
\]

Solution:
Step 1: Divide the numerator and the denominator by $x^2$.
\[
\int \frac{x^2+1}{x^4+1} \, dx = \int \frac{1 + \frac{1}{x^2}}{x^2 + \frac{1}{x^2}} \, dx
\]

Step 2: Apply a algebraic transformation to the denominator.
Recall that $x^2 + \frac{1}{x^2} = \left(x - \frac{1}{x}\right)^2 + 2$.
Substituting this into the integral yields:
\[
\int \frac{1 + \frac{1}{x^2}}{\left(x - \frac{1}{x}\right)^2 + 2} \, dx
\]

Step 3: Solve the integral using substitution.
Let $u = x - \frac{1}{x}$. Differentiating both sides:
\[
du = \left(1 + \frac{1}{x^2}\right) dx
\]
Substituting these values back into the integral:
\[
\int \frac{du}{u^2 + 2}
\]
Using the standard integration formula $\int \frac{dw}{w^2+a^2} = \frac{1}{a} \tan^{-1}\left(\frac{w}{a}\right) + c$:
\[
= \frac{1}{\sqrt{2}} \tan^{-1}\left(\frac{u}{\sqrt{2}}\right) + c
\]

Step 4: Substitute back $u = x - \frac{1}{x}$.
\[
= \frac{1}{\sqrt{2}} \tan^{-1}\left(\frac{x - \frac{1}{x}}{\sqrt{2}}\right) + c = \frac{1}{\sqrt{2}} \tan^{-1}\left(\frac{x^2-1}{\sqrt{2}x}\right) + c
\]

Final Answer:
\[
\boxed{\frac{1}{\sqrt{2}} \tan^{-1}\left(\frac{x^2-1}{\sqrt{2}x}\right) + c}
\]
%%%%%%%%%%%%%%%%%%%%%%%%%%%%%%%%%%%%%%%%%%%%%%%%%%
% QUESTION_END
%%%%%%%%%%%%%%%%%%%%%%%%%%%%%%%%%%%%%%%%%%%%%%%%%%

%%%%%%%%%%%%%%%%%%%%%%%%%%%%%%%%%%%%%%%%%%%%%%%%%%
% QUESTION_START
%%%%%%%%%%%%%%%%%%%%%%%%%%%%%%%%%%%%%%%%%%%%%%%%%%
% id=cse17-final-q5aii
% discipline=CSE
% year=2017
% exam_type=Term Final
% section=B
% ct_number=UNKNOWN
% question_number=5aii
% topic=Indefinite Integration
% subtopic=Trigonometric Integrals
% difficulty=Hard
% length=Long
% question_type=New
% frequency=1
% appearances=
% OCR_STATUS=VERIFIED
\section*{Term Final (Sec B) - Q5(a)(ii)}
Question:
Integrate the following:
\[
\int \frac{dx}{\sec x + \csc x}
\]

Solution:
Step 1: Simplify the integrand in terms of sine and cosine.
\[
\int \frac{dx}{\frac{1}{\cos x} + \frac{1}{\sin x}} = \int \frac{\sin x \cos x}{\sin x + \cos x} \, dx
\]

Step 2: Represent the numerator in terms of the denominator.
Recall that $(\sin x + \cos x)^2 = \sin^2 x + \cos^2 x + 2\sin x \cos x = 1 + 2\sin x \cos x$.
Thus, we can write:
\[
\sin x \cos x = \frac{(\sin x + \cos x)^2 - 1}{2}
\]
Substituting this identity back into the integral:
\[
\int \frac{\frac{1}{2}\left[(\sin x + \cos x)^2 - 1\right]}{\sin x + \cos x} \, dx = \frac{1}{2} \int \left( \sin x + \cos x - \frac{1}{\sin x + \cos x} \right) dx
\]

Step 3: Integrate the terms.
- For the first portion:
  \[
  \frac{1}{2} \int (\sin x + \cos x) \, dx = \frac{1}{2} (\sin x - \cos x)
  \]
- For the second portion, simplify $\sin x + \cos x$ using trigonometric identities:
  \[
  \sin x + \cos x = \sqrt{2} \left( \frac{1}{\sqrt{2}}\sin x + \frac{1}{\sqrt{2}}\cos x \right) = \sqrt{2} \sin\left(x + \frac{\pi}{4}\right)
  \]
  Thus:
  \[
  \frac{1}{2} \int \frac{dx}{\sin x + \cos x} = \frac{1}{2\sqrt{2}} \int \csc\left(x + \frac{\pi}{4}\right) \, dx
  \]
  Using the standard formula $\int \csc w \, dw = \ln|\csc w - \cot w| = \ln\left|\tan\left(\frac{w}{2}\right)\right|$:
  \[
  = \frac{1}{2\sqrt{2}} \ln\left| \tan\left( \frac{x}{2} + \frac{\pi}{8} \right) \right|
  \]

Step 4: Combine the portions.
\[
= \frac{1}{2}(\sin x - \cos x) - \frac{1}{2\sqrt{2}} \ln\left| \tan\left( \frac{x}{2} + \frac{\pi}{8} \right) \right| + c
\]

Final Answer:
\[
\boxed{\frac{1}{2}(\sin x - \cos x) - \frac{1}{2\sqrt{2}} \ln\left| \tan\left( \frac{x}{2} + \frac{\pi}{8} \right) \right| + c}
\]
%%%%%%%%%%%%%%%%%%%%%%%%%%%%%%%%%%%%%%%%%%%%%%%%%%
% QUESTION_END
%%%%%%%%%%%%%%%%%%%%%%%%%%%%%%%%%%%%%%%%%%%%%%%%%%

%%%%%%%%%%%%%%%%%%%%%%%%%%%%%%%%%%%%%%%%%%%%%%%%%%
% QUESTION_START
%%%%%%%%%%%%%%%%%%%%%%%%%%%%%%%%%%%%%%%%%%%%%%%%%%
% id=cse17-final-q5aiii
% discipline=CSE
% year=2017
% exam_type=Term Final
% section=B
% ct_number=UNKNOWN
% question_number=5aiii
% topic=Indefinite Integration
% subtopic=General
% difficulty=Easy
% length=Short
% question_type=New
% frequency=1
% appearances=
% OCR_STATUS=VERIFIED
\section*{Term Final (Sec B) - Q5(a)(iii)}
Question:
Integrate the following:
\[
\int \frac{dx}{\sqrt{x^2 - 4x + 3}}
\]

Solution:
Step 1: Complete the square for the quadratic expression inside the radical.
\[
x^2 - 4x + 3 = (x-2)^2 - 4 + 3 = (x-2)^2 - 1
\]
Substituting this back into the integral:
\[
\int \frac{dx}{\sqrt{(x-2)^2 - 1}}
\]

Step 2: Solve using the standard integration formula.
Recall the standard integration formula:
\[
\int \frac{du}{\sqrt{u^2 - a^2}} = \ln\left| u + \sqrt{u^2 - a^2} \right| + c
\]
Applying this with $u = x-2$ and $a = 1$:
\[
= \ln\left| x - 2 + \sqrt{(x-2)^2 - 1} \right| + c
\]
Simplify the expression inside the radical:
\[
= \ln\left| x - 2 + \sqrt{x^2 - 4x + 3} \right| + c
\]

Final Answer:
\[
\boxed{\ln\left| x - 2 + \sqrt{x^2 - 4x + 3} \right| + c}
\]
%%%%%%%%%%%%%%%%%%%%%%%%%%%%%%%%%%%%%%%%%%%%%%%%%%
% QUESTION_END
%%%%%%%%%%%%%%%%%%%%%%%%%%%%%%%%%%%%%%%%%%%%%%%%%%

%%%%%%%%%%%%%%%%%%%%%%%%%%%%%%%%%%%%%%%%%%%%%%%%%%
% QUESTION_START
%%%%%%%%%%%%%%%%%%%%%%%%%%%%%%%%%%%%%%%%%%%%%%%%%%
% id=cse17-final-q5bi
% discipline=CSE
% year=2017
% exam_type=Term Final
% section=B
% ct_number=UNKNOWN
% question_number=5bi
% topic=Indefinite Integration
% subtopic=General
% difficulty=Medium
% length=Medium
% question_type=New
% frequency=1
% appearances=
% OCR_STATUS=VERIFIED
\section*{Term Final (Sec B) - Q5(b)(i)}
Question:
Show that:
\[
\int \frac{x+\sin x}{1+\cos x} \, dx = x \tan \frac{x}{2} + c
\]

Solution:
Step 1: Substitute trigonometric half-angle identities.
Recall that $1 + \cos x = 2\cos^2\left(\frac{x}{2}\right)$ and $\sin x = 2\sin\left(\frac{x}{2}\right)\cos\left(\frac{x}{2}\right)$.
Substitute these into the integral:
\[
\int \frac{x+\sin x}{1+\cos x} \, dx = \int \frac{x + 2\sin\left(\frac{x}{2}\right)\cos\left(\frac{x}{2}\right)}{2\cos^2\left(\frac{x}{2}\right)} \, dx
\]

Step 2: Separate the integrand into two terms.
\[
= \int \frac{x}{2\cos^2\left(\frac{x}{2}\right)} \, dx + \int \frac{2\sin\left(\frac{x}{2}\right)\cos\left(\frac{x}{2}\right)}{2\cos^2\left(\frac{x}{2}\right)} \, dx
\]
\[
= \frac{1}{2} \int x \sec^2\left(\frac{x}{2}\right) \, dx + \int \tan\left(\frac{x}{2}\right) \, dx
\]

Step 3: Integrate the first term using integration by parts.
For $\int x \sec^2\left(\frac{x}{2}\right) \, dx$, let:
- $u = x \implies du = dx$
- $dv = \sec^2\left(\frac{x}{2}\right) dx \implies v = 2\tan\left(\frac{x}{2}\right)$

Applying the integration by parts formula $\int u \, dv = uv - \int v \, du$:
\[
\int x \sec^2\left(\frac{x}{2}\right) \, dx = 2x \tan\left(\frac{x}{2}\right) - \int 2\tan\left(\frac{x}{2}\right) \, dx
\]

Step 4: Combine the results back into the equation.
\[
\int \frac{x+\sin x}{1+\cos x} \, dx = \frac{1}{2} \left[ 2x\tan\left(\frac{x}{2}\right) - 2 \int \tan\left(\frac{x}{2}\right) \, dx \right] + \int \tan\left(\frac{x}{2}\right) \, dx
\]
\[
= x \tan\left(\frac{x}{2}\right) - \int \tan\left(\frac{x}{2}\right) \, dx + \int \tan\left(\frac{x}{2}\right) \, dx = x\tan\left(\frac{x}{2}\right) + c
\]
Hence proved.

Final Answer:
\[
\boxed{x \tan \frac{x}{2} + c}
\]
%%%%%%%%%%%%%%%%%%%%%%%%%%%%%%%%%%%%%%%%%%%%%%%%%%
% QUESTION_END
%%%%%%%%%%%%%%%%%%%%%%%%%%%%%%%%%%%%%%%%%%%%%%%%%%

%%%%%%%%%%%%%%%%%%%%%%%%%%%%%%%%%%%%%%%%%%%%%%%%%%
% QUESTION_START
%%%%%%%%%%%%%%%%%%%%%%%%%%%%%%%%%%%%%%%%%%%%%%%%%%
% id=cse17-final-q5bii
% discipline=CSE
% year=2017
% exam_type=Term Final
% section=B
% ct_number=UNKNOWN
% question_number=5bii
% topic=Indefinite Integration
% subtopic=General
% difficulty=Medium
% length=Medium
% question_type=New
% frequency=1
% appearances=
% OCR_STATUS=VERIFIED
\section*{Term Final (Sec B) - Q5(b)(ii)}
Question:
Show that:
\[
\int \frac{x^2+1}{(x+1)^2} e^x \, dx = e^x \left(\frac{x-1}{x+1}\right) + c
\]

Solution:
Step 1: Simplify the algebraic fraction.
\[
\frac{x^2+1}{(x+1)^2} = \frac{x^2 - 1 + 2}{(x+1)^2} = \frac{(x-1)(x+1) + 2}{(x+1)^2} = \frac{x-1}{x+1} + \frac{2}{(x+1)^2}
\]
Substituting this back into the integral:
\[
\int e^x \left[ \frac{x-1}{x+1} + \frac{2}{(x+1)^2} \right] dx
\]

Step 2: Identify the pattern of the integrand.
Recall the standard integration rule:
\[
\int e^x [f(x) + f'(x)] \, dx = e^x f(x) + c
\]
Let $f(x) = \frac{x-1}{x+1}$. Differentiating this using the quotient rule:
\[
f'(x) = \frac{1(x+1) - (x-1)(1)}{(x+1)^2} = \frac{2}{(x+1)^2}
\]
This perfectly matches our simplified expression.

Step 3: Evaluate the integral.
\[
\int e^x [f(x) + f'(x)] \, dx = e^x \left(\frac{x-1}{x+1}\right) + c
\]
Hence proved.

Final Answer:
\[
\boxed{e^x \left(\frac{x-1}{x+1}\right) + c}
\]
%%%%%%%%%%%%%%%%%%%%%%%%%%%%%%%%%%%%%%%%%%%%%%%%%%
% QUESTION_END
%%%%%%%%%%%%%%%%%%%%%%%%%%%%%%%%%%%%%%%%%%%%%%%%%%

%%%%%%%%%%%%%%%%%%%%%%%%%%%%%%%%%%%%%%%%%%%%%%%%%%
% QUESTION_START
%%%%%%%%%%%%%%%%%%%%%%%%%%%%%%%%%%%%%%%%%%%%%%%%%%
% id=cse17-final-q6a
% discipline=CSE
% year=2017
% exam_type=Term Final
% section=B
% ct_number=UNKNOWN
% question_number=6a
% topic=Definite Integration
% subtopic=General
% difficulty=Hard
% length=Long
% question_type=New
% frequency=1
% appearances=
% OCR_STATUS=VERIFIED
\section*{Term Final (Sec B) - Q6(a)}
Question:
Evaluate:
\[
\int_0^{\pi/2} \frac{\sin^2 x}{\sin x + \cos x} \, dx
\]

Solution:
Step 1: Use the definite integration property to find an alternative form.
Let $I = \int_0^{\pi/2} \frac{\sin^2 x}{\sin x + \cos x} \, dx$.
Recall the property $\int_0^a f(x) \, dx = \int_0^a f(a-x) \, dx$:
\[
I = \int_0^{\pi/2} \frac{\sin^2\left(\frac{\pi}{2}-x\right)}{\sin\left(\frac{\pi}{2}-x\right) + \cos\left(\frac{\pi}{2}-x\right)} \, dx = \int_0^{\pi/2} \frac{\cos^2 x}{\cos x + \sin x} \, dx
\]

Step 2: Add both equations of $I$ together.
\[
2I = \int_0^{\pi/2} \frac{\sin^2 x + \cos^2 x}{\sin x + \cos x} \, dx = \int_0^{\pi/2} \frac{dx}{\sin x + \cos x}
\]

Step 3: Evaluate the combined integral.
Recall the trigonometric simplification $\sin x + \cos x = \sqrt{2} \sin\left(x + \frac{\pi}{4}\right)$:
\[
2I = \frac{1}{\sqrt{2}} \int_0^{\pi/2} \csc\left(x + \frac{\pi}{4}\right) \, dx
\]
\[
2I = \frac{1}{\sqrt{2}} \left[ \ln\left| \tan\left( \frac{x}{2} + \frac{\pi}{8} \right) \right| \right]_0^{\pi/2}
\]
Evaluate at the limits:
- At $x = \pi/2$: $\tan\left(\frac{\pi}{4} + \frac{\pi}{8}\right) = \tan\left(\frac{3\pi}{8}\right) = \sqrt{2} + 1$
- At $x = 0$: $\tan\left(\frac{\pi}{8}\right) = \sqrt{2} - 1$
\[
2I = \frac{1}{\sqrt{2}} \left[ \ln(\sqrt{2} + 1) - \ln(\sqrt{2} - 1) \right]
\]

Step 4: Simplify the logarithmic term.
Since $(\sqrt{2}-1)(\sqrt{2}+1) = 1 \implies \sqrt{2}-1 = \frac{1}{\sqrt{2}+1}$:
\[
\ln(\sqrt{2}+1) - \ln(\sqrt{2}-1) = \ln(\sqrt{2}+1) - \ln\left(\frac{1}{\sqrt{2}+1}\right) = 2\ln(\sqrt{2}+1)
\]
Substitute this back:
\[
2I = \frac{1}{\sqrt{2}} \cdot 2\ln(\sqrt{2} + 1) \implies I = \frac{1}{\sqrt{2}}\ln(\sqrt{2}+1)
\]

Final Answer:
\[
\boxed{\frac{1}{\sqrt{2}}\ln(\sqrt{2}+1)}
\]
%%%%%%%%%%%%%%%%%%%%%%%%%%%%%%%%%%%%%%%%%%%%%%%%%%
% QUESTION_END
%%%%%%%%%%%%%%%%%%%%%%%%%%%%%%%%%%%%%%%%%%%%%%%%%%

%%%%%%%%%%%%%%%%%%%%%%%%%%%%%%%%%%%%%%%%%%%%%%%%%%
% QUESTION_START
%%%%%%%%%%%%%%%%%%%%%%%%%%%%%%%%%%%%%%%%%%%%%%%%%%
% id=cse17-final-q6b
% discipline=CSE
% year=2017
% exam_type=Term Final
% section=B
% ct_number=UNKNOWN
% question_number=6b
% topic=Definite Integration
% subtopic=General
% difficulty=Hard
% length=Medium
% question_type=New
% frequency=1
% appearances=
% OCR_STATUS=VERIFIED
\section*{Term Final (Sec B) - Q6(b)}
Question:
Show that:
\[
\int_0^{\pi/4} \log(1 + \tan\theta) \, d\theta = \frac{\pi}{8} \log 2
\]

Solution:
Step 1: Set up the integral with the definite integral property.
Let $I = \int_0^{\pi/4} \log(1 + \tan\theta) \, d\theta$.
Using the property $\int_0^a f(x) \, dx = \int_0^a f(a-x) \, dx$:
\[
I = \int_0^{\pi/4} \log\left( 1 + \tan\left( \frac{\pi}{4} - \theta \right) \right) d\theta
\]

Step 2: Simplify the tangent component.
Recall the subtraction formula: $\tan\left(\frac{\pi}{4} - \theta\right) = \frac{1-\tan\theta}{1+\tan\theta}$.
\[
1 + \tan\left(\frac{\pi}{4} - \theta\right) = 1 + \frac{1-\tan\theta}{1+\tan\theta} = \frac{1+\tan\theta + 1-\tan\theta}{1+\tan\theta} = \frac{2}{1+\tan\theta}
\]
Substitute this back:
\[
I = \int_0^{\pi/4} \log\left( \frac{2}{1+\tan\theta} \right) d\theta
\]

Step 3: Separate the logarithmic term.
\[
I = \int_0^{\pi/4} \left[ \log 2 - \log(1+\tan\theta) \right] d\theta
\]
\[
I = \int_0^{\pi/4} \log 2 \, d\theta - I
\]
\[
2I = \log 2 \cdot [\theta]_0^{\pi/4} = \frac{\pi}{4} \log 2
\]
\[
I = \frac{\pi}{8} \log 2
\]
Hence proved.

Final Answer:
\[
\boxed{\frac{\pi}{8} \log 2}
\]
%%%%%%%%%%%%%%%%%%%%%%%%%%%%%%%%%%%%%%%%%%%%%%%%%%
% QUESTION_END
%%%%%%%%%%%%%%%%%%%%%%%%%%%%%%%%%%%%%%%%%%%%%%%%%%

%%%%%%%%%%%%%%%%%%%%%%%%%%%%%%%%%%%%%%%%%%%%%%%%%%
% QUESTION_START
%%%%%%%%%%%%%%%%%%%%%%%%%%%%%%%%%%%%%%%%%%%%%%%%%%
% id=cse17-final-q7a
% discipline=CSE
% year=2017
% exam_type=Term Final
% section=B
% ct_number=UNKNOWN
% question_number=7a
% topic=Definite Integration
% subtopic=General
% difficulty=Hard
% length=Long
% question_type=Repeated PYQ Pattern
% frequency=1
% appearances=
% OCR_STATUS=VERIFIED
\section*{Term Final (Sec B) - Q7(a)}
Question:
Prove that:
\[
\beta(m, n) = \frac{\Gamma(m)\Gamma(n)}{\Gamma(m+n)}
\]

Solution:
Step 1: Write Gamma functions using polar-coordinate substitution.
Recall the definition of the Gamma function:
\[
\Gamma(m) = \int_0^\infty e^{-t} t^{m-1} \, dt
\]
Letting $t = x^2 \implies dt = 2x \, dx$:
\[
\Gamma(m) = 2 \int_0^\infty e^{-x^2} x^{2m-1} \, dx
\]
Similarly:
\[
\Gamma(n) = 2 \int_0^\infty e^{-y^2} y^{2n-1} \, dy
\]

Step 2: Multiply the two Gamma functions together.
\[
\Gamma(m)\Gamma(n) = 4 \int_0^\infty \int_0^\infty e^{-(x^2+y^2)} x^{2m-1} y^{2n-1} \, dx \, dy
\]

Step 3: Convert the double integral into polar coordinates.
Let $x = r\cos\theta$ and $y = r\sin\theta$. The Jacobian is $dx \, dy = r \, dr \, d\theta$.
The boundaries for the first quadrant are $r \in [0, \infty)$ and $\theta \in [0, \pi/2]$:
\[
\Gamma(m)\Gamma(n) = 4 \int_0^{\pi/2} \int_0^\infty e^{-r^2} (r\cos\theta)^{2m-1} (r\sin\theta)^{2n-1} r \, dr \, d\theta
\]
\[
= 4 \int_0^{\pi/2} \int_0^\infty e^{-r^2} r^{2(m+n)-1} \cos^{2m-1}\theta \sin^{2n-1}\theta \, dr \, d\theta
\]
Separate the terms:
\[
= 2 \left( \int_0^\infty e^{-r^2} r^{2(m+n)-1} \, dr \right) \cdot 2 \left( \int_0^{\pi/2} \sin^{2n-1}\theta \cos^{2m-1}\theta \, d\theta \right)
\]

Step 4: Relate the components back to Gamma and Beta definitions.
- The first part is: $\Gamma(m+n)$
- The second part is the trigonometric form of the Beta function: $B(n, m) = B(m, n)$
Therefore:
\[
\Gamma(m)\Gamma(n) = \Gamma(m+n) B(m, n)
\]
\[
B(m, n) = \frac{\Gamma(m)\Gamma(n)}{\Gamma(m+n)}
\]
Hence proved.

Final Answer:
\[
\boxed{B(m, n) = \frac{\Gamma(m)\Gamma(n)}{\Gamma(m+n)}}
\]
%%%%%%%%%%%%%%%%%%%%%%%%%%%%%%%%%%%%%%%%%%%%%%%%%%
% QUESTION_END
%%%%%%%%%%%%%%%%%%%%%%%%%%%%%%%%%%%%%%%%%%%%%%%%%%

%%%%%%%%%%%%%%%%%%%%%%%%%%%%%%%%%%%%%%%%%%%%%%%%%%
% QUESTION_START
%%%%%%%%%%%%%%%%%%%%%%%%%%%%%%%%%%%%%%%%%%%%%%%%%%
% id=cse17-final-q7b
% discipline=CSE
% year=2017
% exam_type=Term Final
% section=B
% ct_number=UNKNOWN
% question_number=7b
% topic=Definite Integration
% subtopic=General
% difficulty=Hard
% length=Medium
% question_type=New
% frequency=1
% appearances=
% OCR_STATUS=VERIFIED
\section*{Term Final (Sec B) - Q7(b)}
Question:
If $U_n = \int_0^1 x^n \tan^{-1} x \, dx$, then prove that:
\[
(n+1)U_n + (n-1)U_{n-2} = \frac{\pi}{2} - \frac{1}{n}
\]

Solution:
Step 1: Integrate $U_n$ by parts.
For $U_n = \int_0^1 x^n \tan^{-1} x \, dx$, let:
- $u = \tan^{-1} x \implies du = \frac{dx}{1+x^2}$
- $dv = x^n \, dx \implies v = \frac{x^{n+1}}{n+1}$

Applying integration by parts:
\[
U_n = \left[ \frac{x^{n+1}}{n+1} \tan^{-1} x \right]_0^1 - \int_0^1 \frac{x^{n+1}}{(n+1)(1+x^2)} \, dx
\]
Multiply both sides by $(n+1)$:
\[
(n+1)U_n = \tan^{-1}(1) - \int_0^1 \frac{x^{n+1}}{1+x^2} \, dx
\]
\[
(n+1)U_n = \frac{\pi}{4} - \int_0^1 \frac{x^{n+1}}{1+x^2} \, dx \quad \text{--- (Equation 1)}
\]

Step 2: Express $(n-1)U_{n-2}$ using a similar form.
By substituting $n$ with $n-2$ in Equation 1:
\[
(n-1)U_{n-2} = \frac{\pi}{4} - \int_0^1 \frac{x^{n-1}}{1+x^2} \, dx \quad \text{--- (Equation 2)}
\]

Step 3: Sum Equation 1 and Equation 2.
\[
(n+1)U_n + (n-1)U_{n-2} = \frac{\pi}{2} - \int_0^1 \frac{x^{n+1} + x^{n-1}}{1+x^2} \, dx
\]

Step 4: Evaluate the remaining integral.
Simplify the integrand:
\[
\frac{x^{n+1} + x^{n-1}}{1+x^2} = \frac{x^{n-1}(x^2+1)}{1+x^2} = x^{n-1}
\]
Thus:
\[
\int_0^1 x^{n-1} \, dx = \left[ \frac{x^n}{n} \right]_0^1 = \frac{1}{n}
\]
Substituting this back yields:
\[
(n+1)U_n + (n-1)U_{n-2} = \frac{\pi}{2} - \frac{1}{n}
\]
Hence proved.

Final Answer:
\[
\boxed{(n+1)U_n + (n-1)U_{n-2} = \frac{\pi}{2} - \frac{1}{n}}
\]
%%%%%%%%%%%%%%%%%%%%%%%%%%%%%%%%%%%%%%%%%%%%%%%%%%
% QUESTION_END
%%%%%%%%%%%%%%%%%%%%%%%%%%%%%%%%%%%%%%%%%%%%%%%%%%

%%%%%%%%%%%%%%%%%%%%%%%%%%%%%%%%%%%%%%%%%%%%%%%%%%
% QUESTION_START
%%%%%%%%%%%%%%%%%%%%%%%%%%%%%%%%%%%%%%%%%%%%%%%%%%
% id=cse17-final-q8a
% discipline=CSE
% year=2017
% exam_type=Term Final
% section=B
% ct_number=UNKNOWN
% question_number=8a
% topic=Applications of Derivatives
% subtopic=General
% difficulty=Medium
% length=Medium
% question_type=New
% frequency=1
% appearances=
% OCR_STATUS=VERIFIED
\section*{Term Final (Sec B) - Q8(a)}
Question:
Find the area of the segment cut off from the parabola $y^2 = 2x$ by the straight line $y = 4x - 1$.

Solution:
Step 1: Find the points of intersection.
From the linear equation:
\[
x = \frac{y+1}{4}
\]
Substitute this into the parabola equation $y^2 = 2x$:
\[
y^2 = 2\left(\frac{y+1}{4}\right) \implies y^2 = \frac{y+1}{2} \implies 2y^2 - y - 1 = 0
\]
Factoring the quadratic equation:
\[
(2y + 1)(y - 1) = 0 \implies y = -\frac{1}{2} \quad \text{or} \quad y = 1
\]

Step 2: Set up the area integration with respect to $y$.
For the interval $y \in [-1/2, 1]$, the line is bounded to the right of the parabola ($x_{line} \ge x_{parabola}$).
The area $A$ is:
\[
A = \int_{-1/2}^1 \left( x_{line} - x_{parabola} \right) dy = \int_{-1/2}^1 \left( \frac{y+1}{4} - \frac{y^2}{2} \right) dy
\]

Step 3: Evaluate the definite integral.
\[
A = \left[ \frac{y^2}{8} + \frac{y}{4} - \frac{y^3}{6} \right]_{-1/2}^1
\]
- At $y = 1$:
  \[
  \frac{1}{8} + \frac{1}{4} - \frac{1}{6} = \frac{3 + 6 - 4}{24} = \frac{5}{24}
  \]
- At $y = -1/2$:
  \[
  \frac{1}{32} - \frac{1}{8} + \frac{1}{48} = \frac{3 - 12 + 2}{96} = -\frac{7}{96}
  \]
Subtracting the lower boundary evaluation:
\[
A = \frac{5}{24} - \left(-\frac{7}{96}\right) = \frac{20}{96} + \frac{7}{96} = \frac{27}{96} = \frac{9}{32}
\]

Final Answer:
\[
\boxed{\frac{9}{32}}
\]
%%%%%%%%%%%%%%%%%%%%%%%%%%%%%%%%%%%%%%%%%%%%%%%%%%
% QUESTION_END
%%%%%%%%%%%%%%%%%%%%%%%%%%%%%%%%%%%%%%%%%%%%%%%%%%

%%%%%%%%%%%%%%%%%%%%%%%%%%%%%%%%%%%%%%%%%%%%%%%%%%
% QUESTION_START
%%%%%%%%%%%%%%%%%%%%%%%%%%%%%%%%%%%%%%%%%%%%%%%%%%
% id=cse17-final-q8b
% discipline=CSE
% year=2017
% exam_type=Term Final
% section=B
% ct_number=UNKNOWN
% question_number=8b
% topic=Applications of Derivatives
% subtopic=General
% difficulty=Medium
% length=Medium
% question_type=New
% frequency=1
% appearances=
% OCR_STATUS=VERIFIED
\section*{Term Final (Sec B) - Q8(b)}
Question:
Find the length of one arc of the cycloid $x = a(\theta - \sin\theta)$, $y = a(1 - \cos\theta)$.

Solution:
Step 1: Set the limits of integration.
One complete arc of a cycloid is generated as the parameter $\theta$ runs from $0$ to $2\pi$.

Step 2: Find the derivatives of $x$ and $y$ with respect to $\theta$.
\[
\frac{dx}{d\theta} = a(1 - \cos\theta)
\]
\[
\frac{dy}{d\theta} = a\sin\theta
\]

Step 3: Setup the arc length differential element.
\[
ds = \sqrt{ \left(\frac{dx}{d\theta}\right)^2 + \left(\frac{dy}{d\theta}\right)^2 } \, d\theta
\]
\[
ds = \sqrt{ a^2(1 - \cos\theta)^2 + a^2\sin^2\theta } \, d\theta
\]
\[
ds = a \sqrt{ 1 - 2\cos\theta + \cos^2\theta + \sin^2\theta } \, d\theta = a\sqrt{2 - 2\cos\theta} \, d\theta
\]
Using the half-angle identity $2 - 2\cos\theta = 4\sin^2\left(\frac{\theta}{2}\right)$:
\[
ds = a \sqrt{4\sin^2\left(\frac{\theta}{2}\right)} \, d\theta = 2a\sin\left(\frac{\theta}{2}\right) d\theta
\]

Step 4: Integrate from $0$ to $2\pi$.
\[
L = \int_0^{2\pi} 2a \sin\left(\frac{\theta}{2}\right) d\theta = \left[ -4a \cos\left(\frac{\theta}{2}\right) \right]_0^{2\pi}
\]
\[
= -4a [\cos\pi - \cos 0] = -4a [-1 - 1] = 8a
\]

Final Answer:
\[
\boxed{8a}
\]
%%%%%%%%%%%%%%%%%%%%%%%%%%%%%%%%%%%%%%%%%%%%%%%%%%
% QUESTION_END
%%%%%%%%%%%%%%%%%%%%%%%%%%%%%%%%%%%%%%%%%%%%%%%%%%

\end{document}